\documentclass[a4paper]{article}
\usepackage{titlesec}
\usepackage{geometry}
\geometry{hmargin={2cm,2cm}, vmargin={3.5cm,3.5cm}}
\usepackage{graphicx}
\usepackage[table]{xcolor}
\usepackage{hyperref}
\hypersetup{pdfborder = {0 0 0},
colorlinks,
citecolor=black,
filecolor=black,
linkcolor=black,
urlcolor=blue}
\definecolor{logo}{RGB}{20,99,163}
\usepackage[T1]{fontenc}
\usepackage[utf8]{inputenc}
\usepackage[english,italian]{babel}
\usepackage{fancyhdr}
\usepackage{lastpage}
\usepackage{comment}

% creazione comando subsubsubsection
\titleclass{\subsubsubsection}{straight}[\subsection]

\newcounter{subsubsubsection}[subsubsection]
\renewcommand\thesubsubsubsection{\thesubsubsection.\arabic{subsubsubsection}}
\titleformat{\subsubsubsection}
    {\normalfont\normalsize\bfseries}{\thesubsubsubsection}{1em}{}
    \titlespacing*{\subsubsubsection}
    {0pt}{3.25ex plus 1ex minus .2ex}{1.5ex plus .2ex}

\makeatletter
\def\toclevel@subsubsubsection{4}
\def\l@subsubsubsection{\@dottedtocline{4}{7em}{4em}}
\makeatother

\setcounter{secnumdepth}{4}
\setcounter{tocdepth}{4}
%fine creazione comando


%Comandi di impaginazione uguale per tutti i documenti
\pagestyle{fancy}
\lhead{\includegraphics[scale=0.2]{../../../../Images/logo.png}}
%Titolo del documento
\rhead{\large Verbale Esterno del 2021-05-21}
%\rfoot{\thepage}
\cfoot{Pagina \thepage{} di \pageref{LastPage}}
\setlength{\headheight}{45pt}
\renewcommand{\footrulewidth}{0.4pt}
\begin{document}

\begin{titlepage}
    \begin{center}
        \includegraphics{../../../../Images/logo.png}\\
        \vspace{20px}
        \textcolor{logo}{\hrulefill}\\
        \vspace{20px}
        \textbf{\huge\textcolor{logo}{Verbale Esterno del 2021-05-21}}\\
        \vspace{10px}
        \textcolor{logo}{\hrulefill}\\
        \vspace{40px}
        \textbf{\Large Informazioni sul documento}\\
        \vspace{20px}
        \begin{tabular}{p{100px} | p{100px}}
            \textbf{Versione}     & 1.0.0                                                                      \\
            \textbf{Responsabile} & Tito Scutari                                                               \\
            \textbf{Redattore}    & Marco Barbaresco                                                           \\
            \textbf{Verificatore} & Amedeo Meggiolaro                                                          \\
            \textbf{Destinatari}  & Prof. Tullio Vardanega \newline Prof. Riccardo Cardin, \newline TechSWEave \\
            \textbf{Stato}        & Approvato                                                                  \\
            \textbf{Uso}          & Esterno                                                                    \\
        \end{tabular}\\
        \vspace{60px}
        \href{mailto:techsweave@gmail.com}{techsweave@gmail.com}\\

    \end{center}
\end{titlepage}

%INDICE
\newpage
\tableofcontents
\newpage

\section{Informazioni riunione}
\begin{itemize}
    \item Data e ora: 2021-05-21 ore 17:30-18:00.
    \item Luogo: Google Meet fornito dai proponenti.
    \item Partecipanti interni: Samuele De Simone, Manuel Varo, Tito Scutari, Nicolò Giaccone, Amedeo Meggiolaro, Marco Barbaresco.
    \item Partecipanti esterni: Milo Ertola, Alessandro Maccagnan.
    \item Segretario: Marco Barbaresco.
\end{itemize}
\section{Argomenti affrontati}
Durante l'incontro sono stati chiesti ai proponenti chiarimenti e consigli riguardo dettagli progettuali e implementativi.\\
I proponenti hanno esaminato alcuni file sorgente, indicandone pregi e difetti.
\section{Verbale della riunione}
\begin{itemize}
    \item Aggiornamento del proponente sullo stato del progetto e sull'esito della \textit{Revisione di Progettazione};
    \item Discussione sulle modalità di comunicazione (servizi Amazon SNS, SQS, SES) tra i vari microservizi;
    \item Discussione sull'indipendenza dei vari microservizi
    \item Esame del codice di alcune funzioni lambda;
    \item Esame dello schema del database;
    \item Discussione sull'importanza della reattività delle pagine del sito;
    \item Discussione sulla struttura degli eventi e su un approccio in grado di rendere più indipendenti i microservizi;
    \item Discussione sul versionamento dei prodotti nel database.
\end{itemize}
\section{Resoconto}
\begin{center}
    \textbf{\Large Registro delle decisioni}\\
    \vspace{10px}
    \begin{table}[h!]
        \centering
        \rowcolors{2}{logo!10}{logo!40}
        \renewcommand{\arraystretch}{1.8}
        \begin{tabular}{p{200px} p{200px}}
            \rowcolor{logo!70} \textbf{Codice} & \textbf{Decisione}                                                                 \\
            VE\_2021\_05\_21.1                 & Utilizzo di Amazon SNS come servizio di messaggistica tra microservizi.            \\
            VE\_2021\_05\_21.2                 & Utilizzo di Amazon SQS come servizio di accodamento messaggi.                      \\
            VE\_2021\_05\_21.3                 & Utilizzo facoltativo e discrezionale di Amazon SES.                                \\
            VE\_2021\_05\_21.4                 & Utilizzo di codice SKU per il versionamento dei prodotti all'interno del database.
        \end{tabular}
    \end{table}
\end{center}

\end{document}